\documentclass[11pt]{article}
\usepackage[T1]{fontenc}
\usepackage{lmodern,amsmath,amssymb,amsthm,mathtools,microtype}
\usepackage[margin=1.15in]{geometry}
\usepackage[hidelinks]{hyperref}
\newtheorem{theorem}{Theorem}[section]
\newtheorem{proposition}[theorem]{Proposition}
\newtheorem{definition}[theorem]{Definition}
\newtheorem{problem}[theorem]{Open Problem}
\title{Compact Hall-to-Borcherds--Kac--Moody Recognition}
\author{Raeez Lorgat}
\date{}
\hypersetup{
  pdftitle={Compact Hall-to-Borcherds--Kac--Moody Recognition},
  pdfauthor={Raeez Lorgat},
  pdfsubject={Finite-height data for Hall algebra recognition}
}
\begin{document}
\maketitle
\begin{abstract}
A denominator records graded multiplicities but forgets the bracket and the
geometric correspondences that should produce it.  We give the elementary
non-recovery obstruction and formulate finite-height data sufficient for a
genuine recognition theorem.  Constructing these data for a compact
Calabi--Yau source remains open.
\end{abstract}

\section{Why the denominator is not the source}

\begin{proposition}[Graded dimensions do not recover a bracket]
There are nonisomorphic positively graded Lie algebras with identical graded
dimensions.
\end{proposition}

\begin{proof}
Let $L$ have basis $x,y$ in degree one and $z$ in degree two, with
$[x,y]=z$ and $z$ central.  Let $L_0$ have the same graded vector space and
zero bracket.  Their graded dimensions agree, while
$[L,L]=\mathbb Cz$ and $[L_0,L_0]=0$, so they are not isomorphic as Lie
algebras.
\end{proof}

Any product depending only on the dimensions of graded root spaces therefore
cannot, by itself, construct the bracket.  A fortiori it does not construct
Hall correspondences, orientation data, a coproduct, or PBW compatibility.

\section{A finite-height recognition datum}

Fix a height bound $N$.  A recognition datum consists of a graded source
bialgebra $H_{\leq N}$ (in graded or super vector spaces), a graded Lie
algebra $\mathfrak b_{\leq N}$, and a graded bialgebra surjection
$\rho_N:H_{\leq N}\twoheadrightarrow U(\mathfrak b_{\leq N})$.  The source
product and coproduct must be constructed from stated correspondences with
the required orientation and pairing data.  The datum must also include a
proof that $\rho_N$ respects the defining relations and a PBW comparison
showing that its asserted kernel, and no larger one, has been imposed.

\begin{problem}
Construct such data from compact Hall correspondences through a nontrivial
height, and construct transition maps in $N$.  For the inverse limit, prove
separatedness, completeness, and the required derived-limit vanishing rather
than replacing a derived limit by a strict one without justification.
\end{problem}
\end{document}
